\documentclass[../main.tex]{subfiles}
\usepackage[utf8]{inputenc}
\usepackage[T1]{fontenc}
\usepackage{graphicx}
\usepackage{longtable}
\usepackage{wrapfig}
\usepackage{rotating}
\usepackage[normalem]{ulem}
\usepackage{amsmath}
\usepackage{amssymb}
\usepackage{capt-of}
\usepackage{hyperref}
\usepackage{float}
\graphicspath{{../}}
\author{Cezary Wieczorkowski}
\date{\today}
\title{Wstęp}
\hypersetup{
 pdfauthor={Cezary Wieczorkowski},
 pdftitle={Wstęp},
 pdfkeywords={},
 pdfsubject={},
 pdflang={Polish}}
\begin{document}


\section{Wstęp i cel pracy}
Sztuczna inteligencja jest jedną z najszybciej rozwijających się dziedzin współczesnej nauki oraz przemysłu. Dzięki jej rozwojowi możliwe stały się rzeczy takie jak generacja realistycznych obrazów na podstawie tekstu lub 
szeroki dostęp do asystentów tekstowych - dużych modeli językowych. Algorytmy sztucznej inteligencji znajdują szerokie zastosowanie również w miejscach gdzie wymagane jest przetwarzanie danych o skomplikowanych do modelowania
zależnościach lub danych charakteryzujących się duża ilością szumu mogącego wynikać z zmiennego środowiska. Przykładem takiej dziedziny są pojazdy autonomiczne które często polegają na algorytmach sztucznej inteligencji 
w celach klasyfikacji danych środowiskowych (np. rozpoznawanie obiektów z kamer) oraz do podejmowania decyzji na podstawie danych z czujników. Istnieje wiele algorytmów sztucznej inteligencji, algorytmem szeroko stosowanym 
w szczególnie w modelach klasyfikujących sygnały oraz obrazy są splotowe sieci neuronowe. 
\par
Pomimo znacznych postępów w dziedzinie sztucznej inteligencji jednym z głównych ograniczeń jej rozwoju pozostaje znaczne zapotrzebowanie 
algorytmów na moc obliczeniową. Duże zapotrzebowanie na moc obliczeniową wynika między innymi z konstrukcji sieci neuronowych. Składają się one z wielu połączonych ze sobą warstw neuronów, w celu uzyskania odpowiedzi sieci neuronowej
na podane pobudzenie należy wykonywać obliczenia kolejno dla każdej warstwy. Obliczenia dla neuronów znajdujących się w jednej warstwie mogą być wykonywane równolegle. Z tego względu sieci algorytmy sieci neuronowych są 
przetwarzane na układach GPU które choć pierwotnie projektowane do obliczeń związanych z grafiką 3D posiadają wiele jednostek obliczeniowych zdolnych do działania równolegle. Podejście to ma niestety duża wadę - układy GPU zużywają
znaczne ilości energii elektrycznej w czasie pracy. Stanowi to dużą przeszkodę przy implementacji algorytmów sztucznej inteligencji w systemach gdzie liczy się oszczędność energii np. pojazdy autonomiczne o zasilaniu bateryjnym.
Jednym z rozwiązań tego problemu jest implementacja dedykowanego akceleratora dla wykorzystywanej sieci w układzie cyfrowym ASIC. Główną wadą tego rozwiązania jest brak możliwości szybkiej oraz taniej możliwości modyfikacji akceleratora.
W związku z tym tracimy jedną z zalet stosowania sieci neuronowej - możliwości ponownego trenowania sieci na zmodyfikowanym zakresie danych w celu rozszerzenia lub poprawy jej funkcjonowania. 
\par
Celem pracy jest zbadanie możliwości implementacji splotowych sieci neuronowych w układach FPGA. Układy te pozwalają na implementacje zróżnicowanych funkcji cyfrowych oraz posiadają możliwość zmiany zaimplementowanej funkcji poprzez
przeprogramowanie układu. Współczesne układy FPGA składają się z dużej ilości bloków logicznych co pozwala na łatwą implementację algorytmów działających równolegle. Te właściwości powodują, że układy tej klasy mają szansę znaleźć
zastosowanie w przypadku implementacji sieci neuronowych w systemach wbudowanych.

\end{document}